% Part: model-theory
% Chapter: models-of-arithmetic
% Section: standard-models

\documentclass[../../../include/open-logic-section]{subfiles}

\begin{document}

\olfileid{mod}{mar}{stm}
\olsection{পাটিগণিতের মানক মডেল}

পাটিগণিতের ভাষা~$\Lang{L_A}$ স্পষ্টতই সংখ্যা নিয়ে, বিশেষ করে
স্বাভাবিক সংখ্যা নিয়ে কথা বলার জন্য অভিপ্রেত। তাই ``সেই'' মানক
মডেল~$\Struct{N}$ বিশেষ: আমরা এই মডেলটি নিয়েই কথা বলতে চাই। কিন্তু
যুক্তিবিদ্যায় প্রায়ই আমাদের আগ্রহ থাকে শুধু গঠনগত বৈশিষ্ট্যে, আর যেকোনো
দুটি সমরূপ !!{structure}s সেই বৈশিষ্ট্যগুলি ভাগ করে। তাই আমরা একটু
উদার হয়ে $\Struct{N}$-এর সমরূপ যেকোনো !!{structure}-কেই ``মানক''
বিবেচনা করতে পারি।

\begin{defn}
$\Lang{L_A}$-এর কোনো !!{structure}, $\Struct{N}$-এর সমরূপ হলে
\emph{মানক}।
\end{defn}

\begin{prop}
\ollabel{prop:standard-domain} কোনো !!a{structure}~$\Struct{M}$ মানক
হলে তার সংজ্ঞাক্ষেত্র মানক সংখ্যাপদগুলির মানের সেট; অর্থাৎ,
\[
\Domain{M} = \Setabs{\Value{\num{n}}{M}}{n \in \Nat}
\]
\end{prop}

\begin{proof}
স্পষ্টতই প্রত্যেক $\Value{\num{n}}{M} \in \Domain{M}$। শুধু দেখাতে হবে,
প্রত্যেক $x \in \Domain{M}$ কোনো~$n$-এর জন্য
$\Value{\num{n}}{M}$-এর সমান। $\Struct{M}$ মানক বলে সেটি
$\Struct{N}$-এর সমরূপ। ধরা যাক, $g\colon \Nat \to \Domain{M}$ একটি
সমরূপতা। তাহলে $g(n) = g(\Value{\num{n}}{N}) =
\Value{\num{n}}{M}$। আবার $g$ !!{surjective} বলে প্রত্যেক
$x \in \Domain{M}$-এর জন্য এমন~$n \in \Nat$ আছে, যাতে $g(n)=x$।
\end{proof}

\begin{explain}
$\Lang{L_A}$-এর কোনো গঠন~$\Struct{M}$ মানক হলে তার !!{domain}-এর সব
উপাদানের নাম মানক সংখ্যাপদ $\num{0}$, $\num{1}$, $\num{2}$, \dots,
অর্থাৎ $\Obj{0}$, $\Obj{0}'$, $\Obj{0}''$ ইত্যাদি পদ দিয়ে দেওয়া যায়।
অবশ্য এর অর্থ এই নয় যে $\Domain{M}$-এর !!{element}s-গুলি
\emph{সত্যিই} সংখ্যা; অর্থ শুধু এই যে~$\Domain{N}$-এর সংখ্যাগুলিকে
যেভাবে নির্দিষ্ট করতে পারি, তাদেরও সেভাবে নির্দিষ্ট করতে পারি।
\end{explain}

\begin{prob}
দেখাও যে \olref[mod][mar][stm]{prop:standard-domain}-এর বিপরীতটি মিথ্যা;
অর্থাৎ এমন একটি !!a{structure}~$\Struct{M}$-এর উদাহরণ দাও, যাতে
$\Domain{M} = \Setabs{\Value{\num{n}}{M}}{n \in \Nat}$, কিন্তু যা
$\Struct{N}$-এর সমরূপ নয়।
\end{prob}

\begin{prop}
\ollabel{prop:thq-standard}
$\Sat{M}{\Th{Q}}$ এবং $\Domain{M} =
\Setabs{\Value{\num{n}}{M}}{n \in \Nat}$ হলে $\Struct{M}$ মানক।
\end{prop}

\begin{proof}
আমাদের দেখাতে হবে, $\Struct{M}$ ও~$\Struct{N}$ সমরূপ। $g(n) =
\Value{\num{n}}{M}$ দিয়ে সংজ্ঞায়িত অপেক্ষক $g\colon \Nat \to
\Domain{M}$ বিবেচনা করো। অনুমান অনুসারে $g$ !!{surjective}। এটি
!!{injective}-ও: $n \neq m$ হলেই $\Th{Q} \Proves
\eq/[\num{n}][\num{m}]$। তাই $\Sat{M}{\Th{Q}}$ হওয়ায়, $n \neq m$
হলেই $\Sat{M}{\eq/[\num{n}][\num{m}]}$। অতএব $n \neq m$ হলে
$\Value{\num{n}}{M} \neq \Value{\num{m}}{M}$, অর্থাৎ $g(n) \neq
g(m)$।

$g$ যে একটি সমরূপতা, তাও যাচাই করতে হবে।
\begin{enumerate}
\item $\Assign{\Obj{0}}{N} = 0$ বলে $g(\Assign{\Obj{0}}{N})=g(0)$।
  $g$-এর সংজ্ঞা অনুসারে $g(0)=\Value{\num{0}}{M}$। কিন্তু $\num{0}$
  ঠিক $\Obj{0}$ পদটি, আর কোনো পদ যদি !!a{constant} হয়, তবে তার মান হলো
  !!{structure}-টি ওই !!{constant}-কে যে মান দেয়; অর্থাৎ
  $\Value{\Obj{0}}{M} =
  \Assign{\Obj{0}}{M}$। অতএব প্রয়োজনমতো
  $g(\Assign{\Obj{0}}{N}) = \Assign{\Obj{0}}{M}$।
\item $\Struct{N}$-এ $\prime$, $\Nat$-এর উপর উত্তরসূরি অপেক্ষক বলে
  $g(\Assign{\prime}{N}(n)) = g(n+1)$। তারপর $g$-এর সংজ্ঞা অনুসারে
  $g(n+1) = \Value{\num{n+1}}{M}$। কিন্তু $\num{n+1}$ ও~$\num{n}'$
  একই পদ, তাই $\Value{\num{n+1}}{M} = \Value{\num{n}'}{M}$। মান
  অপেক্ষকের সংজ্ঞা অনুসারে শেষেরটি
  $= \Assign{\prime}{M}(\Value{\num{n}}{M})$। যেহেতু
  $\Value{\num{n}}{M}=g(n)$, তাই
  $g(\Assign{\prime}{N}(n))=\Assign{\prime}{M}(g(n))$।
\item $\Struct{N}$-এ $+$, $\Nat$-এর উপর যোগ অপেক্ষক বলে
  $g(\Assign{+}{N}(n,m)) = g(n+m)$। তারপর $g$-এর সংজ্ঞা অনুসারে
  $g(n+m)=\Value{\num{n+m}}{M}$। কিন্তু $\Th{Q} \Proves
  \num{n+m} = (\num{n} + \num{m})$, তাই
  $\Value{\num{n+m}}{M}=\Value{\num{n}+\num{m}}{M}$। মান অপেক্ষকের
  সংজ্ঞা অনুসারে শেষেরটি
  $=\Assign{+}{M}(\Value{\num{n}}{M},\Value{\num{m}}{M})$। যেহেতু
  $\Value{\num{n}}{M}=g(n)$ এবং $\Value{\num{m}}{M}=g(m)$, তাই
  $g(\Assign{+}{N}(n, m))=\Assign{+}{M}(g(n), g(m))$।
\item $g(\Assign{\times}{N}(n, m)) = \Assign{\times}{M}(g(n), g(m))$:
  অনুশীলনী।
\item $\tuple{n,m} \in \Assign{<}{N}$ যদি এবং কেবল যদি $n<m$। যদি
  $n<m$ হয়, তবে $\Th{Q} \Proves \num{n} < \num{m}$ এবং
  $\Sat{M}{\num{n}<\num{m}}$-ও হয়। অতএব
  $\tuple{\Value{\num{n}}{M}, \Value{\num{m}}{M}} \in
  \Assign{<}{M}$, অর্থাৎ $\tuple{g(n),g(m)} \in \Assign{<}{M}$। যদি
  $n \not< m$ হয়, তবে $\Th{Q} \Proves \lnot \num{n} < \num{m}$,
  এবং ফলত $\Sat/{M}{\num{n}<\num{m}}$। তাই আগের মতোই
  $\tuple{g(n),g(m)} \notin \Assign{<}{M}$। একত্রে পাই:
  $\tuple{n,m} \in \Assign{<}{N}$ যদি এবং কেবল যদি
  $\tuple{g(n),g(m)} \in \Assign{<}{M}$।
\end{enumerate}
\end{proof}

\begin{explain}
$\Nat$ থেকে $\Lang{L_A}$-এর অন্য যেকোনো !!{structure}~$\Struct{M}$-এর
সংজ্ঞাক্ষেত্রে প্রতিচিত্রণ সংজ্ঞায়িত করার সবচেয়ে স্বাভাবিক উপায় হলো
$g$; কারণ এমন প্রত্যেক $\Struct{M}$-এই $\num{0}$, $\num{1}$,
$\num{2}$ ইত্যাদি দিয়ে নামাঙ্কিত !!{element}s আছে। তাই বিস্ময়ের কিছু
নেই যে, $\Struct{M}$ যদি অন্তত $\num{n}$-গুলি সম্পর্কে কিছু মৌলিক
বিবৃতি $\Struct{N}$-এর মতো একইভাবে সত্য করে এবং $g$-ও একৈক ও সমাপতিত
হয়, তবে $g$ একটি সমরূপতা হয়ে ওঠে। বস্তুত $\Domain{M}$-এ
$\num{n}$-গুলির নামাঙ্কিত !!{element}s ছাড়া আর কিছু না থাকলে এটাই
একমাত্র সমরূপতা।
\end{explain}

\begin{prop}
\ollabel{prop:thq-unique-iso} $\Struct{M}$ মানক হলে
\olref{prop:thq-standard}-এর প্রমাণের $g$-ই $\Struct{N}$ থেকে
$\Struct{M}$-এর উপর একমাত্র সমরূপতা।
\end{prop}

\begin{proof}
ধরা যাক, $h\colon \Nat \to \Domain{M}$ হলো $\Struct{N}$ ও
$\Struct{M}$-এর মধ্যকার একটি সমরূপতা। $n$-এর উপর আরোহ দিয়ে দেখাই যে
$g=h$। যদি $n=0$ হয়, তবে $g$-এর সংজ্ঞা অনুসারে
$g(0)=\Assign{\Obj{0}}{M}$। আবার $h$ সমরূপতা বলে
$h(0)=h(\Assign{\Obj{0}}{N})=\Assign{\Obj{0}}{M}$; অতএব
$g(0)=h(0)$।

এবার $n+1$-এর ক্ষেত্রটি বিবেচনা করি। পাই,
\begin{align*}
  g(n+1) & = \Value{\num{n+1}}{M} \text{, $g$-এর সংজ্ঞা অনুসারে}\\
  & = \Value{\num{n}'}{M} \text{, কারণ $\num{n+1}\ident \num{n}'$}\\
  & = \Assign{\prime}{M}(\Value{\num{n}}{M})
    \text{, $\Value{t'}{M}$-এর সংজ্ঞা অনুসারে}\\
  & = \Assign{\prime}{M}(g(n))  \text{, $g$-এর সংজ্ঞা অনুসারে}\\
  & = \Assign{\prime}{M}(h(n)) \text{, আরোহ-অনুমান অনুসারে}\\
  & = h(\Assign{\prime}{N}(n)) \text{, কারণ $h$ একটি সমরূপতা}\\
  & = h(n+1)
\end{align*}
\end{proof}

\begin{explain}
যেকোনো !!{denumerable} সেট~$M$-এর জন্য $\Nat$ ও~$M$-এর মধ্যে একটি
!!a{bijection} থাকে, তাই এমন প্রত্যেক সেট~$M$ কোনো মানক
মডেল~$\Struct{M}$-এর সম্ভাব্য
!!{domain}। বস্তুত $z \in M$ কোনো বস্তু ও উপযুক্ত কোনো অপেক্ষক~$s$-কে
যথাক্রমে $\Assign{\Obj{0}}{M}$ ও~$\Assign{\prime}{M}$ হিসেবে বেছে
নিলে $+$, $\times$ ও~$<$-এর ব্যাখ্যাও নির্ধারিত হয়ে যায়। মানক মডেলে
$\Assign{\prime}{M}$ হিসেবে কেবল সেই অপেক্ষকগুলি~$s\colon M \to M
\setminus \{z\}$ উপযুক্ত, যেগুলি !!{injective} ও~!!{surjective} উভয়ই।
$s$-এর বিস্তৃতিতে $z$ থাকতে পারে না; থাকলে
$\lforall[x][\eq/[\Obj 0][x']]$ মিথ্যা হতো। ওই !!{sentence}
$\Struct{N}$-এ সত্য,
তাই $\Struct{M}$-কেও সেটি সত্য করতে হবে। অপেক্ষক~$s$-কে !!{injective}
হতে হবে, কারণ $\Struct{N}$-এর উত্তরসূরি অপেক্ষক~$\Assign{\prime}{N}$
তাই, এবং $\Assign{\prime}{N}$ যে !!{injective}, তা $\Struct{N}$-এ সত্য
কোনো !!a{sentence} দিয়ে প্রকাশিত। অপেক্ষকটিকে !!{surjective}-ও হতে হবে;
তা না হলে এমন $x \in M \setminus \{z\}$ থাকত, যা~$s$-এর বিস্তৃতিতে নেই;
অর্থাৎ !!{sentence}~$\lforall[x][(\eq[x][\Obj 0] \lor
\lexists[y][\eq[y'][x]])]$ $\Struct{M}$-এ মিথ্যা হতো---অথচ
$\Struct{N}$-এ সেটি সত্য।
% BN-SRC-121 TARGET-CORRECTION-BEGIN
\par\smallskip\noindent\textnormal{\small\textbf{উৎস-সংশোধন (BN-SRC-121):} হিমায়িত ইংরেজি উৎসের শেষ যুক্তিতে অসমাপতিত~$s$ থেকে কোনো $x$ তার ``domain''-এ নেই বলা হয়েছে। কিন্তু $s\colon M\to M\setminus\{z\}$-এর সংজ্ঞাক্ষেত্র আগেই পুরো~$M$; অসমাপতিত হলে $x$ তার \emph{বিস্তৃতিতে} নেই। বাংলা পাঠে সেই পদটি সংশোধিত।} % BN-SRC-121 TARGET-CORRECTION-END
\end{explain}


\end{document}
